% =============================================================================
\section{Discussion and Limitations}
\label{sec:discussion}
% Budget: 0.5 page. Draft EIGHTH (Step 5.8). FAQ form works well (uncertainty paper).
% At least three honest limitations, each with a mitigation or future-work line.
% The limitations below are the ones the collaborator's handoff says the text must state.
% =============================================================================

\finding{Does this transfer to Linux?} We have not shown it. Linux has the same limit for
anonymous memory: the MMU sets accessed and dirty bits and the kernel samples them. It
already keeps richer versions of the features the learned score uses: \feat{idle}
corresponds to a folio's MGLRU generation age, \feat{sfreq} to MGLRU tiers and DAMON's
access counts, \feat{dirty} to \texttt{PG\_dirty}, and \feat{refaults}/\feat{rdist} to
workingset shadow entries. A learned score could rank the folios of an MGLRU eviction
batch with the same whole-number dot product. The only Linux evidence so far is
indirect: Redis and SQLite traces recorded on Linux with
Valgrind~\citep{nethercote2007valgrind}, replayed in the simulator (\S\ref{sec:eval-info}).

\finding{Why a teaching kernel?} \plan{The whole fault path is small and visible, runs
repeat exactly, and every decision can be traced and replayed, which is what made the
kernel-to-simulator comparison possible. What differs from Linux: per-process (local)
replacement rather than global or cgroup reclaim; one benchmark process at a time on one
hart, so no lock contention on the eviction path; no page cache competing for memory;
timing in QEMU.}

\finding{Limitations.}
\begin{itemize}
\item \textbf{One operating system.} xv6 on QEMU; nothing here says what happens in Linux.
\item \textbf{The fault cost is assumed.} Instruction counts are measured; the time a real
  page-in and write-back take is not, so the break-even results rest on assumed fault
  costs (8\,ms disk, 100\,\textmu s flash) and one instruction per nanosecond.
\item \textbf{Fourteen workloads, some close variants.} 5\% and 20\% memory use one stream
  per variant (51 streams); the intervals show the resulting uncertainty, and most
  differences between the best policies are inside it.
\item \textbf{Adapted classical policies.} ARC, LHD and S3-FIFO ran in the form a kernel
  can run (accessed bits read at evictions), not as published. A per-hit form is
  impossible in a kernel that does not see hits; S3-FIFO's result in particular should not
  be read as the algorithm's.
\item \textbf{Training on six workloads.} ``Global'' means trained on the six original
  workloads; the heavier models ran on only ten streams, so their intervals are the widest.
\item \textbf{No guard in the kernel.} The learned score can be worse than \clock; the
  fallback guard exists only as a simulator prototype.
\item \textbf{Traced data only.} Reference strings cover each workload's data, not its code
  and stack pages (3--9 always-hot frames).
\end{itemize}
% src: week_16 report section "Limits and what is not claimed"; AUSTRA_HANDOFF START HERE sections 1, 4
